\documentclass[11pt,a4paper]{article}
\usepackage[utf8]{inputenc}
\usepackage[T1]{fontenc}
\usepackage{amsfonts}
\usepackage[french]{babel}
\frenchbsetup{StandardLists=true} % à inclure si on utilise \usepackage[francais]{babel}
\usepackage{enumitem}
\usepackage{amssymb}
\usepackage{amsmath}
\usepackage[left=2cm,right=2cm,top=2cm,bottom=2cm]{geometry}
\usepackage{multirow}
\usepackage[dvipsnames]{xcolor}

\usepackage[T1]{fontenc}
\usepackage{fouriernc}
\usepackage{graphicx}

\usepackage{setspace}
\singlespacing

\usepackage{pstricks}
\usepackage{xkeyval}
\usepackage{pst-labo}


\usepackage{multicol}
\usepackage{caption}
\usepackage{fancyhdr}
\pagestyle{fancy}
\renewcommand\headrulewidth{1pt}
\fancyhead[L]{Lycée Denis Diderot }
\fancyhead[R]{Classe de 1TRPM}
\setcounter{page}{2}\fancyfoot[C]{page \thepage}
\fancyfoot[R]{\today}
\fancyfoot[L]{Alexis RUIZ}
\usepackage{multirow,tabularx,booktabs}
\newcommand{\cadreblanc}[1]{%
\medskip\framebox[\linewidth][c]{%
\rule{0mm}{#1mm}}\par
\medskip}

\usepackage{awesomebox}
\setlength{\aweboxvskip}{0mm}
\usepackage{pifont}
\usepackage{chemmacros}
\chemsetup{modules=all}
\setlength{\columnseprule}{.25mm}

\renewcommand{\arraystretch}{2}
\newcommand*\acocher{\quitvmode{\fboxsep0pt \fboxrule0.8pt \fbox{\vrule height1.8ex depth.1ex width0pt \vrule height0pt depth0pt width1.9ex }}}

\newcommand{\code}{\awesomebox[black]{2pt}{

\includegraphics[scale=0.15]{Chapitre 1 - Energie et puissance électrique/images/frame.png}
}{black}}

\newcommand{\moteur}{\awesomebox[black]{2pt}{

\includegraphics[scale=1]{images/moteur.jpg} 
}{black}}

\newcommand{\defconclusion}{\awesomebox[black]{2pt}{\rotatebox{90}{\Large{\textbf{Conclusion}}}}{black}}

\frenchbsetup{StandardLists=true}



\begin{document}

\flushleft \textbf{NOM :\\
Prénom:\\[0,4cm]
}
\hrule\
\\[0,2cm]

\begin{center}
 \LARGE{Activité 1 - \'A quelles grandeurs correspondent les unités électriques ? }
\end{center}

\vspace{3 mm}

\normalsize

\hrule\
\\[0,2cm]
   \begin{center}
   \textbf{Le soin et la rédaction interviendront dans la notation.}
   \end{center} 
   
Une voiture électrique est équipée d'un moteur électrique alimentée par une batterie d'accumulateurs et avec les caractéristiques suivantes : 
 \begin{itemize}
     \item pas d'embrayage. Il suffit d'appuyer sur la pédale d'accélérateur pour que la batterie libère du courant continu. 
     \item pas de boite de vitesses : l'axe du moteur entraîne directement celui des roues motrices. 
     \item un ordinateur embarqué gère, contrôle et régule la puissance à développer. 
     
 \end{itemize}
 
\code{\begin{minipage}{14cm}
{
\begin{itemize} 
    \item Utiliser votre téléphone portable pour scanner le lien puis visualiser la vidéo. 

\end{itemize}  

}
\end{minipage}}
   \vfill
   \ding{202} En vous aidant de la vidéo et du texte ci dessus ; donnez les avantages et les inconvénients du moteur électrique par rapport au moteur thermique. 
\cadreblanc{25}

\vfill 

\noindent%
\begin{minipage}{.4\textwidth}%
\ding{203} Identifiez dans le document ci-contre, les unités électriques dans le Doc.2 et indiquez à quelles grandeurs elle correspondent. \\[3mm]
\cadreblanc{45}

\end{minipage}%
\hfill
\begin{minipage}{.55\textwidth}%
\includegraphics[width=\textwidth]{Chapitre 1 - Energie et puissance électrique/images/A1 voiture.jpg}
\end{minipage}%


\vfill 




\end{document}